\documentclass[11pt,a4paper]{article}
\usepackage[T1]{fontenc}
\usepackage{lmodern}
\usepackage[margin=2.5cm]{geometry}
\usepackage{graphicx}
\usepackage{booktabs}
\usepackage[tableposition=top]{caption}
\usepackage{authblk}
\usepackage[numbers,sort&compress]{natbib}
\usepackage[hidelinks]{hyperref}

\title{Title of the Case: An Unusual Presentation of a Common Condition. A Case Report}
\author[1]{First Author}
\author[2]{Second Author}
\affil[1]{Department of Internal Medicine, University Hospital, City, Country}
\affil[2]{Department of Radiology, University Hospital, City, Country}
\date{}

\begin{document}

\maketitle

\begin{abstract}
\noindent\textbf{Background:} Say what is unusual about this case and why it is worth reporting.\\
\textbf{Case presentation:} Describe the patient's main symptoms, findings, diagnosis, treatment and outcome.\\
\textbf{Conclusions:} Give the main lesson for clinicians in one or two sentences.
\end{abstract}

\noindent\textbf{Keywords:} case report; diagnosis; treatment

\section{Introduction}
Give the background in a few sentences and refer to the relevant literature.
The report follows the CARE guidelines.\cite{gagnier2013care}

\section{Case presentation}
\subsection{Patient information}
Give the patient's age, sex and relevant history, without anything that could
identify them.

\subsection{Clinical findings}
Describe the findings on examination.

\subsection{Timeline}
Table~\ref{tab:timeline} gives the main events in order.

\begin{table}[htbp]
  \centering
  \caption{Timeline of the case.}
  \label{tab:timeline}
  \begin{tabular}{lp{10cm}}
    \toprule
    Day & Event \\
    \midrule
    0 & Presentation to the emergency department with fever and joint pain. \\
    2 & Imaging and laboratory tests; working diagnosis made. \\
    5 & Treatment started. \\
    30 & Follow-up visit: symptoms resolved. \\
    \bottomrule
  \end{tabular}
\end{table}

\subsection{Diagnostic assessment}
Describe the tests, their results (Figure~\ref{fig:imaging}), the differential
diagnosis and how the diagnosis was reached.

\begin{figure}[htbp]
  \centering
  \includegraphics[width=0.5\linewidth]{example-image}
  \caption{Imaging at presentation. Upload an image and use its file name here.}
  \label{fig:imaging}
\end{figure}

\subsection{Therapeutic intervention}
Describe the treatment, its dose and duration, and any changes.

\subsection{Follow-up and outcomes}
Describe the outcome, adherence, and any adverse or unexpected events.

\section{Discussion}
Discuss the strengths and limitations of how the case was managed, the
relevant literature, and what can be learnt from it.

\section{Conclusions}
State the main take-away lesson.

\section*{Patient perspective}
Where possible, the patient's own account of their experience.

\section*{Declarations}
\subsection*{Consent for publication}
Written informed consent was obtained from the patient for publication of this
case report and any accompanying images.

\subsection*{Ethics approval}
State whether the institution requires ethics approval for case reports.

\subsection*{Funding}
State the funding, or that there was none.

\subsection*{Competing interests}
The authors declare no competing interests.

\subsection*{Author contributions}
Describe each author's contribution.

\bibliographystyle{vancouver}
\bibliography{references}

\end{document}
